\chapter{РЕАЛИЗАЦИЯ ПОЛЬЗОВАТЕЛЬСКИХ ПРИЛОЖЕНИЙ}

\section{Требования к пользовательскому интерфейсу}

Пользовательский интерфейс разрабатываемого комплекса ориентирован на инженерные сценарии,
где требуется быстро получать свойства рабочего тела по ограниченному набору входных параметров
и проверять корректность режима по термодинамическим диаграммам.  
Поэтому интерфейс должен поддерживать как одиночные расчёты, так и пакетную обработку наборов точек,
а также визуализацию результатов.  

Базовые функциональные требования к интерфейсу:
\begin{itemize}
\item выбор режима ввода (например, $pT$, $pH$, $pS$, $pX$, $\rho T$) и отображение полного набора свойств;  
\item явное указание единиц измерения и возможность контроля диапазонов;  
\item диагностируемая обработка ошибок и отображение причин отказа в расчёте;  
\item табличный режим с построчной локализацией ошибок;  
\item построение диаграмм (например, $T$--$s$, $p$--$h$, $p$--$v$) и отображение купола насыщения;  
\item экспорт результатов (таблицы, графики, журналы) в распространённые форматы.  
\end{itemize}

Нефункциональные требования:
\begin{itemize}
\item кроссплатформенность и воспроизводимая сборка;  
\item отзывчивость интерфейса при интерактивной работе и при пакетных расчётах;  
\item отсутствие дублирования предметной логики в UI-слое.  
\end{itemize}

Для выполнения этих требований в проекте реализованы две интерфейсные оболочки,
основанные на едином вычислительном ядре \texttt{if97\_core} и общих DTO \texttt{if97\_app\_api}.  
Первая оболочка представляет собой нативное desktop-приложение на FLTK \cite{fltk_rs_docs}.  
Вторая оболочка построена на связке Yew/WASM и Tauri, что обеспечивает современный UI и доступ к функциям ОС \cite{tauri_docs}.  

\section{Настольное приложение на FLTK}

Нативный интерфейс реализован на базе библиотеки FLTK через Rust-биндинги \texttt{fltk-rs} \cite{fltk_rs_docs}.  
Данный вариант ориентирован на локальное использование без web-runtime и сетевой инфраструктуры,
что важно для инженерного рабочего места и демонстрации автономности решения.  

Пользовательская часть организована в виде набора вкладок,
соответствующих типовым сценариям:
«Одиночный расчёт», «Табличный расчёт», «Графики», «О программе».  
Одиночный режим обеспечивает быстрый расчёт одной точки по выбранной входной паре.  
Табличный режим предназначен для пакетной обработки нескольких строк входных данных.  
Графический режим предназначен для визуального анализа режима на диаграммах и для подготовки иллюстративных материалов.  

С архитектурной точки зрения FLTK-клиент реализован как thin-layer над ядром.  
Вся предметная логика остаётся в \texttt{if97\_core}, а UI выполняет ввод, валидацию формата, вызов ядра и представление результата.  
Для повышения надёжности реализации применён подход, близкий к MVU (Model--View--Update),
где обработка событий и изменения состояния происходят через обмен сообщениями.  
Взаимодействие между потоками UI и логикой приложения строится через \texttt{fltk::app::channel},
что позволяет отказаться от «расползающейся» общей мутабельности и уменьшить риск взаимных блокировок.  

Важная особенность FLTK-клиента --- поддержка наборов данных (datasets).  
Это позволяет сохранять и переиспользовать рассчитанные точки и таблицы,
а также управлять их видимостью на диаграммах.  
Такой подход удобен для сравнения нескольких серий расчётов, разных режимов или экспериментальных данных.  

\section{Интерфейс на Yew, WASM и Tauri}

Вторая оболочка реализована на связке web-технологий и desktop-упаковки.  
Frontend построен на фреймворке Yew и компилируется в WebAssembly.  
Desktop-функции, вычисления и доступ к файловой системе реализованы в backend-части Tauri \cite{tauri_docs}.  

Принципиальный момент архитектуры состоит в том, что frontend не выполняет термодинамические расчёты самостоятельно.  
Вместо этого UI формирует запросы в формате DTO (слой \texttt{if97\_app\_api}) и вызывает команды Tauri через механизм \texttt{invoke}.  
Tauri backend, в свою очередь, выполняет вычисления через \texttt{if97\_core} и возвращает сериализуемый результат обратно в UI.  
Это обеспечивает единый источник истины по вычислениям и совпадение результатов с другими каналами доставки.  

Для обмена данными важна совместимость с сериализацией в JSON и типичными ограничениями web-окружения.  
Поэтому DTO-слой учитывает необходимость корректного кодирования специальных числовых значений (например, \texttt{NaN} и \texttt{Infinity})
и хранения диагностических сообщений об ошибках.  

Термодинамические расчёты относятся к CPU-bound задачам.  
Если запускать их внутри асинхронного runtime без изоляции,
возможны зависания и деградация отзывчивости интерфейса.  
Поэтому тяжёлые вычисления в Tauri backend выполняются в блокирующем пуле через \texttt{spawn\_blocking},
что отделяет вычислительную нагрузку от UI и I/O.  

На уровне UI-архитектуры решается задача сохранения введённых данных и результатов при переключении вкладок.  
Для этого применяется принцип «единого источника истины»: состояние вынесено в общий контекст приложения,
а навигация реализована через управление видимостью (CSS-классы),
без разрушения DOM-узлов и потери пользовательского ввода.  

\section{Табличные расчёты, диаграммы, купол насыщения и экспорт данных}

Обе интерфейсные реализации поддерживают два базовых режима работы: одиночный и табличный.  
Одиночный режим используется для интерактивного ввода одной точки и анализа полного набора свойств состояния.  
Табличный режим ориентирован на пакетные инженерные сценарии,
когда требуется обработать массив входных данных и получить итоговую таблицу результатов.  

Ключевое требование табличного режима --- устойчивость к некорректным строкам.  
В реальных данных отдельные строки могут выходить за пределы области применимости IF97 или содержать ошибки ввода.  
Поэтому таблица обрабатывается построчно,
а ошибки локализуются на уровне конкретной строки без остановки всего расчёта.  

Визуализация результатов выполняется через построение термодинамических диаграмм.  
Диаграммы используются для контроля режима и для наглядного представления результатов в отчёте.  
Для корректной интерпретации состояния на диаграмме важна отрисовка купола насыщения,
то есть границ $x=0$ (насыщенная жидкость) и $x=1$ (насыщенный пар).  

Построение купола насыщения выполняется с учётом требований к гладкости и стоимости вычислений.  
Для диаграмм, где по оси откладывается удельный объём $v$,
применяются меры для повышения плотности точек в области малых значений $v$,
в частности использование метрики, основанной на $\ln(v)$.  
В web/Tauri-оболочке дополнительно используется кэширование результатов по типу диаграммы и варианту осей,
а также защита от устаревших асинхронных ответов при частом изменении настроек отображения.  

Экспорт данных необходим как для инженерной работы, так и для подготовки материалов к защите.  
В проекте поддерживаются следующие типовые варианты экспорта:
\begin{itemize}
\item сохранение табличных результатов в CSV;  
\item экспорт диаграмм в растровые изображения (PNG) для вставки в текст ВКР;  
\item выгрузка журналов (logs) и диагностических сообщений для анализа ошибок и воспроизводимости результатов.  
\end{itemize}

В FLTK-версии экспорт реализуется как сохранение подготовленных файлов из нативного приложения.  
В Tauri-версии экспорт выполняется через backend-команды и системные диалоги,
что обеспечивает корректную интеграцию с ОС и единый пользовательский опыт на разных платформах \cite{tauri_docs}.  

\section{Сравнение интерфейсных реализаций и сценариев применения}

Наличие двух независимых оболочек демонстрирует переносимость вычислительного ядра и корректность архитектурного решения «одно ядро --- несколько потребителей».  
Обе реализации используют \texttt{if97\_core}, поэтому при одинаковом вводе обеспечивают совпадение численных результатов.  
Различия относятся к технологическому стеку, способам доставки и эксплуатационным сценариям.  

FLTK-версия является максимально лёгкой и автономной.  
Она подходит для быстрого запуска, офлайн-использования и работы на маломощных ПК,
а также удобна для демонстрации минимального слоя между пользователем и ядром.  

Yew/WASM + Tauri-версия ориентирована на современный UI,
кроссплатформенную упаковку и интеграцию с системными возможностями (файловые диалоги, экспорт, журналы) \cite{tauri_docs}.  
Эта реализация удобна для расширения интерфейса, визуализации и сценариев,
где важны единые UX-паттерны и переносимость web-подходов на desktop.  

Таким образом, две оболочки решают одинаковую предметную задачу,
но оптимальны для разных условий применения,
а единый вычислительный слой обеспечивает корректность и воспроизводимость результатов.  
